\section{The real-place determinant estimates}\label{sec:real-place}

We estimate the determinant defined in Equation~\eqref{eq:determinant}
without a rationality assumption.  Write
\begin{equation}\label{eq:real-normalized-parameters}
 \begin{gathered}
 \alpha=\frac an=\frac{11}{48},\quad
 \beta=\frac bn=\frac7{48},\quad
 \gamma=\frac gn=\frac qn=\frac4{48},\quad
 \eta=\frac hn=\frac2{48},\\
 D_*=n-1-2g=40N-1.
 \end{gathered}
\end{equation}
In particular, $0<D_*<n$ and $n\ge48$.  For a real list $y$ of
length $n$, let
\[
 V(y)=\prod_{i<j}(y_j-y_i),\qquad \mathcal V(y)=|V(y)|.
\]
All constants in the estimates of this Section are independent of the
locations and separations of the integration variables.

\subsection{The integral and its two sheets}

Put $d\mu(t)=|t|\,dt/f(t)$ on $(-1,1)$, and set
\[
 A_r(t)=P_r(t)-\frac32\boldsymbol 1_{\{t>0\}}\frac{f(t)}tD_r(t),
 \qquad \psi_k(s)=s^{b+k}(1-s)^q.
\]
The value at $t=0$ is immaterial; the displayed expression has a finite
limit there because $D_r(t)/t$ is a polynomial.  The determinant entries
are exactly
\[
 \int_{-1}^1\int_0^1 A_r(t)\frac{\psi_k(s)}{1-ts}\,ds\,d\mu(t).
\]
The measure has mass $\mu((-1,1))=2$.  The scalar majorant
$\int\!\int(1-|t|s)^{-1}\,ds\,d\mu(t)$ is finite: integration in $s$
gives at worst a logarithmic singularity at $|t|=1$, which is integrable
against $(1-t^2)^{-1/2}dt$.  The functions $A_r$ and $\psi_k$ are
bounded for each fixed $n$.  Thus all permutation expansions below are
absolutely integrable.

We apply Andr\'eief's determinant integration identity twice
\cite[Equation~(1.7) and Section~2.2]{Forrester2018Andreief}.
Expanding determinants and relabeling integration variables gives
the identity
\begin{equation}\label{eq:real-double-integral}
 \Delta_N=\frac1{(n!)^2}\int_{(-1,1)^n}\int_{(0,1)^n}
 \det[A_r(t_i)]_{r,i}\,
 \det\left[\frac1{1-t_i s_j}\right]_{i,j}\,
 \det[\psi_k(s_j)]_{j,k}\,ds\,d\mu^n(t).
\end{equation}
Here $0\le r,k<n$ and $1\le i,j\le n$.  To see the factor $1/n!$
at each application directly, expand the two determinants sharing an
integration list: every one of the $n!$ permutations of that list gives
the same determinant of single integrals.  Fubini's theorem applies by
the preceding majorant, also after each permutation expansion.

Cauchy's double alternant in multiplicative variables
\cite[Section~2.1, Equation~(2.7)]{Krattenthaler1999} is
\begin{equation}\label{eq:real-cauchy-determinant}
 \det\left[\frac1{1-t_i s_j}\right]_{i,j}
 =\frac{V(t)V(s)}{\prod_{i,j}(1-t_i s_j)}.
\end{equation}
For completeness, multiplying by the denominator gives a polynomial
alternating separately in $t$ and $s$, of degree at most $n-1$ in each
individual variable.  Dividing by $V(t)V(s)$ therefore leaves a constant.
Expanding $(1-t_i s_j)^{-1}$ at $s=0$, the first nonzero homogeneous
part of the determinant is $V(t)V(s)$, so that constant is one.  Also
$\det[\psi_k(s_j)]_{j,k}=V(s)\prod_j s_j^b(1-s_j)^q$.

Use the real coordinate
\begin{equation}\label{eq:real-coordinate}
 x_i=\frac{t_i}{1+\sqrt{1-t_i^2}},\qquad
 t_i=\frac{2x_i}{1+x_i^2},\qquad
 d\mu(t_i)=\frac{4|x_i|}{(1+x_i^2)^2}\,dx_i.
\end{equation}
The coordinate maps $(-1,1)$ bijectively to $(-1,1)$.  Except on a
set of measure zero we may assume that all $x_i$ are distinct and
nonzero.  Define
\[
 (\xi_{\mathrm f}(x),\xi_{\mathrm n}(x))=
 \begin{cases}(1/2,1/2),&x<0,\\(-1/4,5/4),&x>0,\end{cases}
 \qquad
 \mathcal R_n(x)=
 \det\left[\xi_{\mathrm f}(x_i)x_i^{g-r}
             +\xi_{\mathrm n}(x_i)x_i^{r-g}\right]_{r,i}.
\]
The subscripts $\mathrm n$ and $\mathrm f$ refer to the near and far
evaluations at $x_i$ and $1/x_i$, respectively, inside and outside the
unit circle.
Indeed, $fD_r/t=(R_r^*-R_r)/2$, so the negative half-interval
has the row $(R_r+R_r^*)/2$, whereas the positive half-interval has
$(5R_r-R_r^*)/4$.  Factoring $(1-t_i)^ht_i^{C-1}$ out of the
$i$th column therefore gives the exact identity
\begin{equation}\label{eq:real-sheet-integral}
 \Delta_N=\frac1{(n!)^2}\int\!\int
 \mathcal R_n(x)\,
 \frac{V(t)V(s)^2}{\prod_{i,j}(1-t_i s_j)}
 \prod_j s_j^b(1-s_j)^q
 \prod_i t_i^{C-1}(1-t_i)^h\,ds\,d\mu^n(t).
\end{equation}
In particular, there is one $t$ Vandermonde and two $s$
Vandermondes.  All the denominators in this formula are positive.

\subsection{A uniform interpolating function}

The interpolation problem comes from factoring out the far-sheet weight.
Put $c_i=\xi_{\mathrm n}(x_i)/\xi_{\mathrm f}(x_i)$; the denominator
is nonzero on both half-intervals. For $0\le r<n$,
\begin{equation}\label{eq:real-sheet-factorization}
 \xi_{\mathrm f}(x_i)x_i^{g-r}
   +\xi_{\mathrm n}(x_i)x_i^{r-g}
 =\xi_{\mathrm f}(x_i)x_i^{g-(n-1)}
   \left(x_i^{n-1-r}+c_i x_i^{D_*}x_i^r\right).
\end{equation}
Thus a holomorphic function with values $h_*(x_i)=c_i x_i^{D_*}$
turns the expression in parentheses into
$x_i^{n-1-r}+h_*(x_i)x_i^r$. We shall bound the resulting evaluation
determinant by controlling the norm of this function. The first estimate
applies when
\begin{equation}\label{eq:case-condition}
 \sum_{i=1}^n\frac{1-x_i^2}{1+x_i^2}\le D_*.
\end{equation}
The following construction is uniform even when the nodes cluster.

\begin{lemma}\label{lem:real-uniform-interpolation}
Let $x_1,\ldots,x_n$ be distinct nonzero real numbers in $(-1,1)$
satisfying Equation~\eqref{eq:case-condition}.  There is a function
$h_*$ holomorphic on a neighborhood of the closed unit disk such that
\[
 h_*(x_i)=c_i x_i^{D_*},\qquad
 c_i=\begin{cases}1,&x_i<0,\\-5,&x_i>0,\end{cases}
 \qquad
 \sup_{|z|\le1}|h_*(z)|\le K_0 n,
 \quad K_0=10e^{12}.
\]
The neighborhood may depend on the node list; the displayed bound does not.
\end{lemma}

\begin{proof}
Set
\[
 B(z)=\prod_i\frac{z-x_i}{1-x_i z},\qquad
 F(z)=\frac{z^{D_*}}{B(z)}.
\]
For $0<u\le1$ and a real $x$ with $0<|x|<1$, the logarithmic
derivative of one factor on the imaginary diameter is
\[
 \frac{d}{d\log u}\frac12\log\frac{u^2+x^2}{1+x^2u^2}
 =\frac{(1-x^4)u^2}{(u^2+x^2)(1+x^2u^2)}
 \le\frac{1-x^2}{1+x^2}.
\]
The inequality follows on dividing the denominator by $u^2$ and
using $u^2+u^{-2}\ge2$.  Integration from $u$ to $1$ gives
\[
 |B(iu)|\ge u^{\sum_i(1-x_i^2)/(1+x_i^2)}\ge u^{D_*}.
\]
The same holds at $-iu$, and $F(0)=0$.  For $1\le u\le1+2/n$,
each factor has modulus at least one because
$(u^2+x^2)-(1+x^2u^2)=(u^2-1)(1-x^2)\ge0$.
Consequently
\begin{equation}\label{eq:real-diameter-bound}
 |F(iu)|\le e^2\qquad (|u|\le1+2/n).
\end{equation}

Orient the fixed segment $\Gamma=[-i(1+2/n),i(1+2/n)]$ upwards,
and, off that segment, define the fixed Cauchy integral
\[
 \mathcal C(z)=\frac1{2\pi i}\int_\Gamma\frac{6F(w)}{w-z}\,dw.
\]
Write $c_-=1$, $c_+=-5$, with the sign specifying the left or right
half-plane, and put
\begin{equation}\label{eq:real-glued-interpolant}
 h_*(z)=c_\pm z^{D_*}-B(z)\mathcal C(z)
 \quad\text{when }\ \pm\operatorname{Re}z>0.
\end{equation}
The density is holomorphic in a neighborhood of every point of
$\Gamma$, since its poles $x_i$ are nonzero and real.  Local contour
deformation and the Cauchy integral formula show that the left lateral
value of $\mathcal C$ minus its right lateral value is $6F$.
More explicitly, move a short upward piece of $\Gamma$ to its left;
the closed contour formed by the original piece followed by the reversed
new piece is positively oriented, and its residue at $w=z$ is $6F(z)$.
Thus the difference of the two local analytic continuations in
Equation~\eqref{eq:real-glued-interpolant} is
$(c_--c_+)z^{D_*}-6B(z)F(z)=0$.  This proves holomorphic gluing
across the segment, including its two crossings of the unit circle.

For each fixed node list the poles $1/x_i$ of $B$ and the endpoints
of $\Gamma$ lie strictly outside the closed unit disk.  A sufficiently
small exterior neighborhood avoids all these points, and the same local
gluing works there.  It follows that $h_*$ is holomorphic on that
neighborhood.  Since $B(x_i)=0$ and $\mathcal C$ is regular at the
nonzero real point $x_i$, the required interpolation follows.

It remains to bound this one fixed function.  In either radius-$1/4$
disk about $i$ or $-i$, the quantities $|w|$, $|w-x_i|$, and
$|1-x_iw|$ are at least $3/4$.  Hence $F$ is nonzero there and
\begin{equation}\label{eq:real-logarithmic-derivative}
 \left|\frac{F'(w)}{F(w)}\right|
 \le\frac43D_*+\frac83n<4n.
\end{equation}
As $|F(i)|=|F(-i)|=1$, integration along a segment in either disk
gives $|F(w)|\le e^{12}$ whenever $|w-i|\le3/n$ or
$|w+i|\le3/n$.

If $|z|=1$ and $|\operatorname{Re}z|\ge1/(10n)$, the original
segment has distance at least $1/(10n)$ from $z$; its length is at
most $3$.  Equation~\eqref{eq:real-diameter-bound} gives
$|\mathcal C(z)|\le30e^2n$.
If instead $0<|\operatorname{Re}z|<1/(10n)$, $z$ is near one
of $\pm i$.  Near $i$, replace the portion of $\Gamma$ from
$i(1-1/n)$ to $i(1+2/n)$ by the three sides of the rectangle whose
other vertical side has real part
$-\operatorname{sign}(\operatorname{Re}z)/n$.  Near $-i$, make
the identical replacement of the bottom portion with imaginary parts
from $-1-2/n$ to $-1+1/n$.  The swept rectangle is in the opposite
half-plane from $z$, and it is contained in the radius-$1/4$ disk
about the relevant endpoint.  It contains neither $z$ nor a pole of
$F$, so this deformation leaves the value of the fixed integral
$\mathcal C(z)$ unchanged.

Each new side is within $3/n$ of $i$ or $-i$.  The new contour has
length at most $3$ and distance at least $1/(2n)$ from $z$.
For the latter assertion the vertical side has horizontal separation
at least $1/n$, the outer horizontal side has vertical separation
at least $2/n$, and the inner horizontal side and remaining diameter
have separation at least
\[
 \frac1n-\bigl(1-\sqrt{1-(10n)^{-2}}\bigr)>\frac1{2n}.
\]
Thus $|\mathcal C(z)|\le6e^{12}n$ in this case.  Only the two
endpoint neighborhoods were deformed; the middle diameter uses its
absolute bound \eqref{eq:real-diameter-bound}, with no derivative
estimate near zero.  Finally $|B(z)|=1$ on $|z|=1$, so
Equation~\eqref{eq:real-glued-interpolant} gives
$|h_*(z)|\le K_0n$ there.  At $z=\pm i$ use continuity of the
glued function.  The maximum modulus principle proves the bound in
the disk.  No uniform lower bound for the exterior neighborhood was used.
\end{proof}

\subsection{Evaluation in a finite-dimensional Hilbert space}

The following argument uses the kernel and compression viewpoint of
bounded analytic interpolation; see \cite{Sarason1967}.
We prove the needed finite-dimensional facts directly.

\begin{proposition}[Interpolation estimate]\label{prop:first-sheet-bound}
Under Equation~\eqref{eq:case-condition},
\begin{equation}\label{eq:real-first-sheet-bound}
 |\mathcal R_n(x)|\le(1+K_0n)^n
       \mathcal V(1/x)\prod_i|x_i|^g.
\end{equation}
In particular, the prefactor is $\exp(O(n\log n))$ uniformly in the nodes.
\end{proposition}

\begin{proof}
Let $Q(z)=\prod_i(1-x_i z)$, and give
\[
 \mathcal H_Q=\{v/Q:\ v\in\mathbb C[z],\ \deg v<n\}
\]
the norm inherited from the Hilbert space $H^2$ of analytic functions
with square-summable Taylor coefficients:
\[
 \|v/Q\|^2=\frac1{2\pi}\int_0^{2\pi}
          \frac{|v(e^{i\theta})|^2}{|Q(e^{i\theta})|^2}\,d\theta.
\]
For each fixed list the zeros of $Q$ lie outside the closed disk, so
this is a well-defined norm.  Define
$\widetilde v(z)=z^{n-1}v(1/z)$ and
$R(v/Q)=\widetilde v/Q$.  This is a complex-linear involution and
an isometry: on the circle
$|\widetilde v(e^{i\theta})|=|v(e^{-i\theta})|$, while the real
coefficients of $Q$ imply
$|Q(e^{i\theta})|=|Q(e^{-i\theta})|$.

The functions $k_i(z)=(1-x_i z)^{-1}$ belong to $\mathcal H_Q$.
They are linearly independent, as their distinct poles $1/x_i$ show,
and hence form a basis.  For $u\in H^2$, the Taylor coefficient
inner product gives $\langle u,k_i\rangle=u(x_i)$.
If $\Pi_Q$ is orthogonal projection from $H^2$ to $\mathcal H_Q$,
it follows that
\[
 (\Pi_Q u)(x_i)=u(x_i)\quad(1\le i\le n).
\]
Multiplication by the function $h_*$ of
Lemma~\ref{lem:real-uniform-interpolation} has $H^2$ operator norm
at most $\|h_*\|_\infty$.  Consequently
\[
 J=\Pi_Q M_{h_*}|_{\mathcal H_Q},\qquad L=I+JR
 \quad\text{satisfy}\quad \|L\|\le1+K_0n.
\]

Let $E:\mathcal H_Q\to\mathbb C^n$ send $v/Q$ to
$(v(x_i))_i$.  It is invertible because a polynomial of degree less
than $n$ cannot vanish at all the distinct nodes.
For $u\in\mathcal H_Q$, $E(u)_i=Q(x_i)u(x_i)$.
Since $LR=R+J$ and projection preserves function values, applying
this identity to $LR(v/Q)$ cancels the factors $Q(x_i)$
algebraically and gives
\begin{equation}\label{eq:real-evaluation-cancellation}
 E L R(v/Q)=\bigl(\widetilde v(x_i)+h_*(x_i)v(x_i)\bigr)_i.
\end{equation}
In the basis $1/Q,z/Q,\ldots,z^{n-1}/Q$, $\det E=V(x)$,
and $R$ has determinant of absolute value one.  Hence the determinant
of the mixed evaluations divided by $V(x)$ has absolute value
$|\det L|\le\|L\|^n\le(1+K_0n)^n$.  This is an algebraic
cancellation of the evaluation determinant.  In particular no bound
for the inverse evaluation matrix, which might be poorly conditioned,
enters the argument.

Apply this bound to the mixed evaluations in
Equation~\eqref{eq:real-sheet-factorization}. Since
$|\xi_{\mathrm f}(x_i)|\le1/2\le1$ and
$\mathcal V(1/x)=\mathcal V(x)\prod_i|x_i|^{-(n-1)}$,
Equation~\eqref{eq:real-first-sheet-bound} follows.
\end{proof}

\begin{proposition}[Hadamard estimate]\label{prop:second-sheet-bound}
For every distinct nonzero real list in $(-1,1)$, and in particular
when Equation~\eqref{eq:case-condition} fails,
\begin{equation}\label{eq:real-second-sheet-bound}
 |\mathcal R_n(x)|\le
 (3/2)^n n^{n/2}2^{-\binom n2}
 \prod_i(1+x_i^2)^{(n-1)/2}|x_i|^{g-(n-1)}.
\end{equation}
The factor $(3/2)^n n^{n/2}$ is $\exp(O(n\log n))$; the power
$2^{-\binom n2}$ is retained in the principal integrand below.
\end{proposition}

\begin{proof}
Expand the determinant by choosing one sheet in each column.  For
$z_i\in\{x_i,1/x_i\}$ the absolute monomial determinant is
$\mathcal V(z)\prod_i|z_i|^{-g}$.  Write
$z_i=\tan\phi_i$ with $-\pi/2<\phi_i<\pi/2$.  Then
\[
 \mathcal V(z)=\prod_i(1+z_i^2)^{(n-1)/2}
              \prod_{i<j}|\sin(\phi_j-\phi_i)|.
\]
The last product is $2^{-\binom n2}$ times the Vandermonde on
the unit-circle points $e^{2i\phi_i}$.  Its monomial matrix has
column norms $\sqrt n$, so Hadamard's inequality gives the bound
$2^{-\binom n2}n^{n/2}$.  The ratio of the remaining weight on
the far sheet to that on the near sheet is
\[
 \frac{|1/x|^{-g}(1+x^{-2})^{(n-1)/2}}
      {|x|^{-g}(1+x^2)^{(n-1)/2}}
 =|x|^{-D_*}\ge1.
\]
Thus all the weights can be bounded by their far-sheet values.
Finally the sum of absolute sheet coefficients is at most
$\prod_i(|\xi_{\mathrm f}(x_i)|+|\xi_{\mathrm n}(x_i)|)
\le(3/2)^n$, proving the result.
\end{proof}

\subsection{The two principal integrands}

To state precisely the quantities needed for the real-place energy
estimate, put $t_i=2x_i/(1+x_i^2)$ and
\[
 \mathcal J_n(x,s)=
 \frac{\mathcal V(t)\mathcal V(s)^2}{\prod_{i,j}(1-t_i s_j)}
 \prod_j s_j^b(1-s_j)^q
 \prod_i |t_i|^{C-1}(1-t_i)^h.
\]
Define
\begin{equation}\label{eq:real-integrand-majorants}
 \begin{split}
 \mathcal I_{2,n}(x,s)
   &=\mathcal J_n(x,s)\mathcal V(1/x)\prod_i|x_i|^g,\\
 \mathcal I_{1,n}(x,s)
   &=\mathcal J_n(x,s)2^{-\binom n2}
       \prod_i(1+x_i^2)^{(n-1)/2}|x_i|^{g-(n-1)}.
 \end{split}
\end{equation}
Let $\Omega_{2,n}$ be the set of interior configurations with distinct
nonzero $x_i$ satisfying Equation~\eqref{eq:case-condition}, and let
$\Omega_{1,n}$ be its complementary case among such configurations.
We use the interpolation estimate on $\Omega_{2,n}$ and the
Hadamard estimate on $\Omega_{1,n}$; the latter estimate is valid
in both cases. After substituting $t=2x/(1+x^2)$, $\kappa$ is the
exponent of the absolute $x$-Vandermonde in the corresponding
majorant: $\kappa=2$ for the interpolation estimate and
$\kappa=1$ for the Hadamard estimate.
The $s_j$ always range over $(0,1)$.

The preceding propositions and Equation~\eqref{eq:real-sheet-integral}
give
\[
 |\Delta_N|\le\frac{2^n K_n}{(n!)^2}
 \max_{\kappa\in\{1,2\}}\sup_{(x,s)\in\Omega_{\kappa,n}}
      \mathcal I_{\kappa,n}(x,s),
 \quad
 K_n=\max\{(1+K_0n)^n,(3/2)^n n^{n/2}\}.
\]
Indeed $d\mu^n(t)\,ds$ has total mass $2^n$, and the two case
sets partition its domain up to a null set.  In logarithmic form,
with $\log0=-\infty$, this yields
\begin{equation}\label{eq:real-supremum-reduction}
 \frac{\log|\Delta_N|}{n^2}-\frac12\log2
 \le\max_{\kappa\in\{1,2\}}
 \sup_{(x,s)\in\Omega_{\kappa,n}}
 \left(\frac{\log\mathcal I_{\kappa,n}(x,s)}{n^2}
                -\frac12\log2\right)
 +O\left(\frac{\log(n+2)}n\right).
\end{equation}
The error is independent of both integration lists.  Factorials,
measure masses, and sheet sums have all been accounted for explicitly.

Repeated nonzero $x_i$ give a zero original row determinant and are
also handled by continuity in the first estimate; repeated $s_j$
give a zero $s$ Vandermonde.  The sets with a zero node or a boundary
node have measure zero for the absolutely continuous measures above.
At such points the original expression in
Equation~\eqref{eq:real-double-integral}, rather than its factored
Laurent expression, supplies the integrable definition.  Approaches
to these exceptional sets require no separation condition in any
bound proved here.  In particular all configurations arbitrarily
close to them remain included in the suprema in
Equation~\eqref{eq:real-supremum-reduction}.
